\chapter*{下编\quad 考释}
\phantomsection
\addcontentsline{toc}{chapter}{下编\quad 考释}
\markboth{下编\quad 考释}{下编\quad 考释}
\chaptersubtitle{ЧАСТЬ ВТОРАЯ}{ИССЛЕДОВАНИЕ}
\setcounter{chapter}{0}

\noindent
本编八章。每章从一片或几片残简出发，说明它能证明什么，不能证明什么。凡上编已经解决的问题，此处不再重复；凡上编留下的空白，此处也不假装能填上。

八章之间不是并列的。第一章论人，第二章论这套人维持自身的方式，第三章论这套方式的一个特例，第四章论这套方式的核心机关，第五章回到数字，第六章论结束，第七章论与传世文献的关系，第八章论名号。读者若只读一章，建议读第五章。

\chapter{月之民：一群改过自己的人}

简十四把两个数字并排放着：计口三十一万二千，实居十一万九千。这两句话说的不是一回事，但它们挨在一起，说明写它的人认为它们该挨在一起。

这两个数之间差十二万一千。上编卷二对此只字未解。本章要说明的是，这个差额不是账目上的疏漏，是一群人对自己的处置留下的记录。

\kao{一\quad 他们不是这里的土著}

月之民来自何处，本卷说不清楚。可以确定的只有一点：他们到得比这些字更早，而他们到的时候，这里还什么都没有。

证据在简一。「以石为垣，垣厚四丈，高十二丈」，「地无草木，风自西来，终岁不息」。这是一处新建之地的记录：垣的尺寸、风的方向、草木的缺席，全是要写下来的，因为一切都是新的。

故土不需要记这些。住在故土的人早就知道风从哪边来，不需要刻在金属上留档。

简一记的是营建，不是迁徙。换言之，这三十七片开口的时候，事情已经过去很久了。

\kao{二\quad 他们把自己改过}

简十三：「新入者，先居外舍三月，乃入环。入环之日，给籍，给衣，给器。」

这一条的对象是自地上归来的人，但它反过来证明一件事：所谓「入环」，不只是进一道门，是进入一个条件不同之处。人要先在外舍住满三个月，才能进去。

三个月里做什么，本卷无一字。但简四、简五给了侧面：气路、水路、外工、浴水，「出则浴，浴水不入渠」，「衣别贮，不与内同」。一个人进出，要换气、换水、沐浴、换衣，还要在垣外住满三个月，然后才能进去。

这些安排的密度，与他处所记的日常生活不成比例。一个只需要防尘的地方不需要这样。

\kao{三\quad 改造的代价}

他们改了什么，《战争的雷霆》说得很清楚，此处只需要与本卷对上。

那一套改造的结果，是他们不再变化：不老化，不再生育，也不再适应新的环境。作为交换，他们活得极长，长到能把一座城维持两万年。

这个交换是自愿的。简十三的「给籍，给衣，给器」说明入环是一道手续，不是一种惩罚。没有人被赶进去。

简十五「不足则减，减自上始」也说明这一点。减自上始，先减居者，后减工，后减匠。若城是一个牢笼，次序会反过来：先减最不要紧的人。此处的次序说明，在这个体系里，居者最低，匠最高，而这套等级是所有人都认的。

\kao{四\quad 关于「不会死」}

须把这一条分清，否则本书与《战争的雷霆》会打架。

月之民并非不死。他们会死于外伤，会死于极大量的污秽，也会因为处在一个变化不断的地方而缓慢地老去、死去。那一本书说他们会死，只是死得慢；本书与此不冲突。

简二十六的四个存活率数字，就是「死得慢」的具体样子。四代之内，从六一四降到四二，中间没有一次骤降，全是缓慢的滑落。

「四代之后，无记」五个字下面还有五个字：「无记，非无也」。这五个字是后人补的，字口与上文不同。补它的人显然明白，空白本身会被误读成不存在。

\kao{五\quad 器上的名字}

有一处细节与全书的气氛相反，须记在这里。

简十三：「器不继，以旧器补。旧器有记，记其旧主。」

这是一个以「名可削」为制度的地方（简二十二第十七项），却在器物上留着人的名字。撤离时削名，日常里记名，两件事并存。

我不认为这里有矛盾。削名是撤离时的处置，记名是日常的惯例；一个是终结，一个是延续。但两件事同时存在于一个体系里，说明这个体系对「名」的看法并不统一。

这一条留到第八章再说。

\chapter{污秽：变化本身}

简五有一句话，读起来平平无奇：「出则浴，浴水不入渠。」

若不看全书，这只是一条卫生规定。看了全书，它是一整套制度的入口。

\kao{一\quad 干净是维持出来的}

月之都的干净不是天生的。这一点必须强调，因为几乎所有关于「月」的传说都假定那里本来就干净。

简四：「气自下而上，水自上而下，皆以管。管在壁中，不可见。凡管，岁一检。」

这句话描述的不是一座天然的洁净之城，是一套需要每年检修的系统。天然干净的地方不需要检修，更不需要「岁一检」这种制度化的检修。

检修还有一套更细的安排：「检者不入环，于道中为之。检毕，录其数于板，板悬于门。」

检修者不进环，在环外的道上作业；检完把数记在板上，板挂在门口。这是把一个技术动作变成一次公开的记录。若只是修管子，不必如此。

\kao{二\quad 污秽是什么}

本卷没有一条残简给出定义。简五只说了处置，简十九只说了数目，简二十二第十一项只说「山，封」。

但有一条间接的线索：所有这些制度，都指向同一件事，即把变化挡在外面，或者把已经进来的变化弄出去。

气要过滤，水要循环，人要沐浴，外工要住在垣外，垣外百步不得置物，焚籍之后的灰要埋到垣外。这一整套安排针对的不是灰尘，也不是细菌，是「变」。

在这个体系里，「脏」与「变」是同一件事，「干净」与「不变」是同一件事。我找不到反例。

\kao{三\quad 三重处置}

若「污秽」即「变」，则这套制度应当同时管三样东西：物、人、言。

物有处置，见简六、简十八。人有处置，见简五、简十三。言有处置，见简九：「山名不可呼。呼者其罪与越界同。」

第三样最可注意。禁止一个名字被说出，与禁止一个人越界，判为同罪。把「说」与「走」等量齐观，说明在这个体系里，言语不是言语，是行为的一种。

通行本《老子》有「名可名，非常名」的话，与此形似而实不同。那一句说的是名不能尽物，这一条说的是名不能出诸口。前者是认识论，后者是治安。

\kao{四\quad 界}

简九是全书最难读的一片：

\begin{center}
\small
界者，非垣非墙，非水非火。在其内者出，必失其半。失者不复。
\end{center}

「界」是什么，本卷无一字解释。它不是垣，不是墙，不是水，不是火；这四句用的是排除法，而排除法不告诉你它是什么。

「必失其半」是什么意思，也无从判定。或指所携之物，或指其人。「失者不复」则说明这种减损不可逆。

我不打算为它给一个说法。它就在那里，是全书最大的一处空白，比「四代之后无记」更彻底，因为那一处至少知道缺的是什么。

\kao{五\quad 一个反例}

若要给这一章找一个反例，简十三的「旧器有记，记其旧主」是最合适的一条。

旧器上有旧主的名字。按本章的读法，名是可以被说出、被削去、被禁止的东西；可是在这里，名字被刻在器物上，随器物一起流转，一再被补给下一个人。

这与简九「山名不可呼」正好相反：一个名字连说都不能说，另一个名字却堂而皇之地留在器上。

两件事如何并存，本卷不能回答。我把它列在这里，与其说是反例，不如说是这一章读法的一个界限。

\chapter{纯化：分解为纯粹之物}

上编三十七片里，与纯化可能有关的只有半句：简二十二第十三项，「药，弃」。

全书没有第二处。本章要说明的是，这半句为什么值得单列一章，以及为什么它撑不起一章。

\kao{一\quad 一个「弃」字}

「药，弃」的「药」，在《战争的雷霆》里有专名，叫蓬莱药。那一本书说它产自净秽山，由月兔加工。本卷没有一条残简印证这条生产链，只有这一处「弃」。

一个「弃」字能说明的事，比一整条生产线更多：到撤离的时候，这种东西不再需要了。

撤离清单十九项，处置方式分四类：封、断、毁、弃。药归入「弃」，与「毁」不同。「毁」是把东西弄坏，「弃」是把它扔掉。弄坏的不能再用，扔掉的还在某处。

扔在哪里，本片未言。

\kao{二\quad 捣与拆}

按《战争的雷霆》的引述，纯化是一种能力，其作用是把有形状的东西拆开，拆回没有形状的状态。这个说法在残简里找不到对应。

简十六「捣者三千」、简十七「人日三合」说的都是「捣」，与「拆」未必是一回事。捣是把东西弄碎，拆是把东西分开；弄碎之后仍是碎块，分开之后各部分还在。二者的产物不同。

我不认为这两者可以互相印证。本卷的材料不够，硬要对上，只会把两件事都弄坏。

\kao{三\quad 这是生产，不是仪式}

可以确定的只有一点：这套体系里存在一种「把东西变成更简单的状态」的操作，且有定额、有计量、有专职工人。

简十七：「人日三合，多者赏，少者罚。罚不过三，过三则移籍。赏以粟，罚以役。」

一个工序有定额，有超额奖励，有不足惩罚，有退出机制，还有与之配套的结算方式。这不是仪式，是生产。仪式不做定额，也不设退出。

三千人日捣三合，一日九千合。这个量级说明它不是为少数人准备的。

\kao{四\quad 纯狐不在此卷}

「纯狐」二字在三十七片中一次也没有出现。

这个名字见于后世文献，与月之都的关系是本丛书另一卷的题目。本章不作推测。

同样，「嫦娥」「后羿」也不在此卷。三十七片里没有一个人名，这是本书与《战争的雷霆》最大的差别：那一本处处是人，这一本一个名字也没有。

简二十二第十七项「名，削」，或许正是这个现象的原因；但那是撤离时的事，解释不了两万年。

\chapter{净秽山与月兔}

简九说「界外有山，山名不可呼」。简二十二第十一项说「山，封」。

两处说的应当是同一座山。本卷对它的记载仅此，其余全部来自后世的推断。

\kao{一\quad 材料分布的不均衡}

这是全书材料分布最不均衡的一处。

城本身有九片残简记其营建，人口有十片记其岁计，撤离有十二片记其清点，而这座山，按《战争的雷霆》的说法是整座城能够维持下去的关键，在本卷里只有六个字。

我不会为了让它显得重要而多写。它的重要性不是从本卷看出来的，是从城的存续时间看出来的：一座两万年不用换气的城，必然有一个足够大的处理段。

顺着这条线索推算：十一万九千人，日给粟二升，即二千三百八十石；人日三合之捣，即九千合。这两组数字与「山」之间有什么关系，本卷无一句交代。

\kao{二\quad 只封不毁}

十九项退役措施之中，十项是封，四项是断，三项是毁，一项是弃，一项是移。

山排在第十一项，处置是「封」。城可以毁，路可以断，器可以毁，药可以弃，唯独山只封。

封与毁的差别，在于可否恢复。毁是不可逆的，封是可逆的：封起来的东西将来还能打开。撤离者给山留了一条后路，给城没有留。

\kao{三\quad 月兔两行}

月兔在本卷里只占两行。

简十六「捣者三千」，是唯一一处可能指月兔的记载。但「捣者」只是一个工种，本卷没有说明他们是谁、从哪来、与他处的人有什么不同。

简十九「兔乱一万一千」，是第二处。这场乱事在《战争的雷霆》里叫月兔之乱，西方文献作「月面战争」。

两项相加，月兔在本卷里造成了一万一千人死亡，然后就不再出现。既没有记载他们是什么，也没有记载乱后如何处置。

\kao{四\quad 那八千}

有一件事必须在这里说清楚，否则读者会被卷二的数字骗过去。

简十九的三项加起来是二万二千：兔乱一万一千，他故八千，逃亡三千。而简十四与简二十之间的差额是三万。

差八千。

这八千人在残简里没有任何一条记载。它们既不在事故之列，也不在撤离之列，也不在留存之列。

下一章会回到这个数字，并说明它可能有三种去处。
